
% !TeX root = ../main.tex

\section{Einleitung und Problemstellung}\label{sec:einleitung}

Die zuverlässige Authentifizierung von Benutzern und
technischen Systemen stellt eine zentrale Herausforderung
des Identity and Access Managements (IAM) dar.
Bevor einem Benutzer der Zugriff auf schützenswerte
Ressourcen gewährt werden kann, muss zunächst überprüft
werden, ob dieser tatsächlich über die erforderlichen
Authentifizierungsmerkmale verfügt.

Klassische Authentifizierungsverfahren basieren häufig
auf Passwörtern, die als gemeinsames Geheimnis zwischen
einem Benutzer und einem authentifizierenden System
dienen. Solche Verfahren können jedoch verschiedene
Sicherheitsrisiken mit sich bringen. Dazu zählen
insbesondere schwache oder wiederverwendete Passwörter,
Phishing sowie die Kompromittierung serverseitig
gespeicherter Passwort-Hashes.

Asymmetrische kryptografische Verfahren verfolgen
einen anderen Ansatz. Anstelle eines gemeinsamen
Geheimnisses wird ein Schlüsselpaar verwendet, das
aus einem privaten und einem öffentlichen Schlüssel
besteht. Während der private Schlüssel ausschließlich
dem jeweiligen Schlüsselinhaber zur Verfügung stehen
soll, kann der öffentliche Schlüssel anderen
Teilnehmern zugänglich gemacht werden.

Hieraus ergibt sich eine grundlegende kryptografische
Problemstellung: Wie kann eine Partei nachweisen,
dass sie den zu einem öffentlichen Schlüssel
gehörenden geheimen Schlüssel kennt, ohne diesen
gegenüber einer anderen Partei offenzulegen?


\subsection{Problemstellung}

Für die Betrachtung dieser Problemstellung werden
zwei beteiligte Parteien unterschieden:

\begin{description}

    \item[Prover (beweisende Partei):]
    Der Prover besitzt einen geheimen Schlüssel und
    möchte nachweisen, dass er diesen kennt.

    \item[Verifier (prüfende Partei):]
    Der Verifier kennt den zugehörigen öffentlichen
    Schlüssel und möchte überprüfen, ob der Prover
    tatsächlich über das entsprechende geheime
    Wissen verfügt.

\end{description}

Eine naheliegende, jedoch unsichere Möglichkeit
bestünde darin, den geheimen Schlüssel unmittelbar
an den Verifier zu übertragen. Dieser könnte
anschließend überprüfen, ob der übermittelte
Schlüssel zum bekannten öffentlichen Schlüssel
gehört.

Ein solcher Ansatz würde allerdings die Geheimhaltung
des privaten Schlüssels aufheben. Der Verifier
oder ein Angreifer, der Zugriff auf den übertragenen
Schlüssel erlangt, könnte diesen für zukünftige
Authentifizierungsvorgänge verwenden.

Ein geeignetes kryptografisches Verfahren muss
daher einen Nachweis der Schlüsselkenntnis
ermöglichen, ohne dass der geheime Schlüssel
selbst übertragen werden muss.

Darüber hinaus sollte ein aufgezeichneter
Authentifizierungsvorgang nicht ohne Weiteres
für eine spätere Authentifizierung wiederverwendet
werden können.

Für die praktische Anwendung ist außerdem eine
vertrauenswürdige Zuordnung zwischen dem
öffentlichen Schlüssel und der jeweiligen
Benutzeridentität erforderlich.

Der mathematische Nachweis der Kenntnis eines
privaten Schlüssels bestätigt zunächst lediglich,
dass eine Partei über das entsprechende geheime
Wissen verfügt. Ob dieser Schlüssel tatsächlich
einer bestimmten Person zugeordnet werden kann,
hängt von der zuvor erfolgten Registrierung
und Identitätsprüfung ab.

Die kryptografische Identifikation ist somit
ein Bestandteil eines umfassenderen
Authentifizierungs- und Identitätsmanagementsystems.


\subsection{Das Schnorr-Identifikationsprotokoll}

Eine Lösung für die beschriebene Problemstellung
stellt das Schnorr-Identifikationsprotokoll dar.

Das Verfahren geht auf den deutschen Mathematiker
und Kryptografen Claus-Peter Schnorr zurück.
In seiner Veröffentlichung \emph{Efficient Signature
Generation by Smart Cards} aus dem Jahr 1991
beschreibt Schnorr effiziente Verfahren zur
kryptografischen Identifikation und zur Erzeugung
digitaler Signaturen \cite{schnorr1991}.

Die mathematische Grundlage des Verfahrens bildet
das diskrete Logarithmusproblem in geeigneten
zyklischen Gruppen.

Dabei wird aus einem geheimen Exponenten $x$
und einem öffentlich bekannten Gruppenelement $g$
der öffentliche Schlüssel
\[
    y=g^x
\]
berechnet.

Während die Berechnung des öffentlichen Schlüssels
aus dem geheimen Exponenten effizient möglich ist,
gilt die umgekehrte Berechnung des geheimen
Exponenten aus dem öffentlichen Schlüssel bei
geeigneter Wahl der mathematischen Gruppe und
ihrer Parameter als rechnerisch schwierig.

Das Schnorr-Identifikationsprotokoll nutzt diese
Eigenschaft, um einen interaktiven Nachweis der
Kenntnis des geheimen Exponenten zu ermöglichen.

Der Nachweis erfolgt durch den Austausch von
drei aufeinanderfolgenden Nachrichten:

\begin{enumerate}
    \item Commitment,
    \item Challenge,
    \item Response.
\end{enumerate}

Der Verifier überprüft anhand einer mathematischen
Gleichung, ob die empfangene Response mit dem
öffentlichen Schlüssel und den zuvor ausgetauschten
Werten übereinstimmt.

Der geheime Schlüssel selbst wird während der
Protokollausführung nicht übertragen.

Eine besondere Eigenschaft des Schnorr-Protokolls
besteht darin, dass sich die Kenntnis des
geheimen Schlüssels nachweisen lässt, ohne
einem ehrlichen Verifier zusätzliche Informationen
über diesen Schlüssel offenzulegen.

Diese Eigenschaft wird als Honest-Verifier
Zero Knowledge bezeichnet und im weiteren
Verlauf der Arbeit näher untersucht.


\subsection{Abgrenzung zu anderen kryptografischen Verfahren}

Das Schnorr-Identifikationsprotokoll muss von
anderen Verfahren der asymmetrischen Kryptografie
unterschieden werden.

Insbesondere verfolgen Verschlüsselung, digitale
Signaturen und kryptografische Identifikation
unterschiedliche Schutzziele.

\paragraph{Verschlüsselung.}

Das primäre Ziel einer Verschlüsselung besteht
darin, die Vertraulichkeit von Informationen
sicherzustellen.

Eine Nachricht wird mithilfe eines kryptografischen
Verfahrens so transformiert, dass ihr ursprünglicher
Inhalt nur von berechtigten Parteien wiederhergestellt
werden kann.

Das Schnorr-Identifikationsprotokoll dient hingegen
nicht der Verschlüsselung von Nachrichten.
Es ermöglicht den Nachweis kryptografischen Wissens,
ohne die Vertraulichkeit der übertragenen
Kommunikationsinhalte selbst sicherzustellen.

\paragraph{Digitale Signaturen.}

Digitale Signaturen ermöglichen die Überprüfung
der Authentizität und Integrität einer Nachricht.

Hierzu erzeugt der Inhaber eines privaten Schlüssels
eine Signatur für eine bestimmte Nachricht.
Andere Teilnehmer können diese Signatur mithilfe
des zugehörigen öffentlichen Schlüssels überprüfen.

Im Unterschied dazu dient das ursprüngliche
Schnorr-Identifikationsprotokoll dem interaktiven
Nachweis der Kenntnis eines geheimen Schlüssels
gegenüber einem Verifier.

Eine unmittelbare Bindung des Nachweises an eine
zu signierende Nachricht ist im ursprünglichen
Identifikationsprotokoll nicht vorgesehen.

Allerdings lassen sich aus dem Schnorr-Protokoll
durch geeignete Erweiterungen digitale
Signaturverfahren ableiten.

Insbesondere die Fiat-Shamir-Transformation
ermöglicht es, die interaktive Challenge durch
einen kryptografischen Hashwert zu ersetzen.
Dieser Zusammenhang bildet eine Grundlage
für Schnorr-Signaturen.

\paragraph{Kryptografische Identifikation.}

Das Ziel eines kryptografischen
Identifikationsverfahrens besteht darin,
einen Teilnehmer anhand eines ihm zugeordneten
kryptografischen Geheimnisses zu authentifizieren.

Das Schnorr-Protokoll verfolgt dieses Ziel,
indem es die Kenntnis eines geheimen
diskreten Logarithmus nachweist.

Die Identifikation einer konkreten Person
setzt dabei weiterhin voraus, dass der
verwendete öffentliche Schlüssel zuvor
vertrauenswürdig mit ihrer Identität
verknüpft wurde.

Tabelle~\ref{tab:kryptografische-verfahren}
fasst die unterschiedlichen Zielsetzungen
zusammen.

\begin{table}[htbp]
    \centering
    \small
    \renewcommand{\arraystretch}{1.3}
    \begin{tabular}{p{3.5cm}p{8cm}}
        \toprule
        Verfahren & Primäres Ziel\\
        \midrule

        Verschlüsselung
            & Schutz der Vertraulichkeit
              einer Nachricht\\

        Digitale Signatur
            & Nachweis der Authentizität
              und Integrität einer Nachricht\\

        Schnorr-Identifikation
            & Nachweis der Kenntnis eines
              geheimen Schlüssels ohne dessen
              Offenlegung\\

        \bottomrule
    \end{tabular}
    \caption{Abgrenzung ausgewählter
    kryptografischer Verfahren.
    Eigene Darstellung.}
    \label{tab:kryptografische-verfahren}
\end{table}


\subsection{Zielsetzung und Forschungsfrage}

Ziel dieser Arbeit ist es, die mathematischen
Grundlagen und die Funktionsweise des
Schnorr-Identifikationsprotokolls zu untersuchen
sowie dessen Sicherheit und praktische
Einsetzbarkeit im Kontext des Identity and
Access Managements zu bewerten.

Dabei soll insbesondere nachvollzogen werden,
wie ein Teilnehmer die Kenntnis eines
geheimen Schlüssels nachweisen kann, ohne
diesen gegenüber der prüfenden Partei
offenzulegen.

Aus dieser Zielsetzung ergibt sich folgende
zentrale Forschungsfrage:

\begin{quote}
    \textit{Wie ermöglicht das
    Schnorr-Identifikationsprotokoll den Nachweis
    der Kenntnis eines geheimen Schlüssels ohne
    dessen Offenlegung, und unter welchen
    mathematischen und technischen Voraussetzungen
    eignet sich das Verfahren für eine sichere
    und effiziente Authentifizierung?}
\end{quote}

Zur Beantwortung dieser Forschungsfrage werden
folgende Teilfragen untersucht:

\begin{enumerate}

    \item Auf welchen mathematischen Grundlagen
    basiert das Schnorr-Identifikationsprotokoll?

    \item Wie erfolgt der Nachweis der Kenntnis
    eines geheimen Schlüssels innerhalb des
    Drei-Schritt-Protokolls?

    \item Welche kryptografischen
    Sicherheitseigenschaften besitzt das Verfahren,
    und unter welchen Voraussetzungen gelten diese?

    \item Welcher rechnerische Aufwand ist für
    die reguläre Protokollausführung und für
    mögliche Angriffe erforderlich?

    \item Welche Vorteile, Nachteile und
    praktischen Einsatzbedingungen ergeben
    sich für die Verwendung des Verfahrens
    in Authentifizierungssystemen?

\end{enumerate}

Die Untersuchung erfolgt anhand einer theoretischen
Betrachtung der mathematischen Grundlagen,
der Analyse des Protokollablaufs sowie einer
literaturgestützten Bewertung seiner
Sicherheitseigenschaften.

Zur Veranschaulichung werden numerische Beispiele
mit kleinen Gruppenparametern verwendet.
Diese dienen ausschließlich der Nachvollziehbarkeit
der mathematischen Zusammenhänge und stellen
keine kryptografisch sicheren Parametersätze dar.

Ergänzend werden theoretische Aufwandsabschätzungen
für ausgewählte Angriffsverfahren vorgenommen.

Eine experimentelle Untersuchung konkreter
Implementierungen oder eine vollständige
Sicherheitsanalyse sämtlicher möglicher
Angreifermodelle ist nicht Gegenstand dieser
Arbeit.


\subsection{Aufbau der Arbeit}

Die vorliegende Arbeit gliedert sich in
sechs Abschnitte.

Nach der Einführung in die Problemstellung
und der Formulierung der Forschungsfrage
werden in Abschnitt~\ref{sec:grundlagen}
die mathematischen Grundlagen erläutert.

Hierzu zählen insbesondere Gruppen,
zyklische Gruppen, modulare Arithmetik
und das diskrete Logarithmusproblem.

Abschnitt~\ref{sec:protokoll} beschreibt
anschließend das Schnorr-Identifikationsprotokoll.

Neben der Schlüsselerzeugung und dem Ablauf
des Drei-Schritt-Protokolls werden dessen
Korrektheit sowie die wesentlichen
kryptografischen Sicherheitseigenschaften
untersucht.

In Abschnitt~\ref{sec:praxis} erfolgt
eine Betrachtung des praktischen Bezugs
und des rechnerischen Aufwands.

Dabei werden insbesondere generische
Verfahren zur Berechnung diskreter
Logarithmen, theoretische Angriffszeiten
und ausgewählte Anwendungsfälle betrachtet.

Abschnitt~\ref{sec:bewertung} untersucht
die Vor- und Nachteile des Verfahrens
und leitet daraus geeignete
Einsatzbedingungen ab.

Hierzu wird das Schnorr-Protokoll auch
mit anderen Authentifizierungsmechanismen
verglichen.

Abschließend werden in
Abschnitt~\ref{sec:fazit} die gewonnenen
Erkenntnisse zusammengeführt, die
Forschungsfrage beantwortet und die
Grenzen der vorliegenden Untersuchung
aufgezeigt.
